
% ============================================================
% download-gitbash.tex — A SUBDOCUMENT
% Compiles standalone ("pdflatex git-bash-windows.tex") using
% main.tex's preamble, OR gets inserted into main.tex via
% \subfile{git-bash-windows}.
% ============================================================

\documentclass[../main.tex]{subfiles}

\begin{document}

\section{Downloading and Installing Git Bash on Windows}
% Top-level heading — becomes a numbered section if pulled
% into main.tex.

\subsection{What is Git Bash?}
Git Bash is a terminal application for Windows that lets you run
Git commands (and many Unix-style commands like \texttt{ls}, \texttt{cd},
\texttt{ssh-keygen}) using a Bash shell, instead of the default
Windows Command Prompt.

\subsection{Download the Installer}
\begin{enumerate}[leftmargin=*]
    \item Open a web browser and go to \url{https://git-scm.com/downloads}.
    % \url{...} (from the hyperref package) makes the link
    % clickable in the PDF and displays it in monospace font.
    \item Click \textbf{Windows} under the list of operating systems.
    \item The download should start automatically. If not, click the
    manual download link for the latest 64-bit version.
\end{enumerate}

\subsection{Run the Installer}
\begin{enumerate}[leftmargin=*]
    \item Locate the downloaded file (usually named something like
    \texttt{Git-2.xx.x-64-bit.exe}) in your \textbf{Downloads} folder
    and double-click it.
    \item If prompted by Windows ``Do you want to allow this app to
    make changes to your device?'', click \textbf{Yes}.
    \item Click \textbf{Next} on the license screen.
\end{enumerate}

\subsection{Installation Options}
The installer presents several configuration screens. For most
beginners, the default options are fine — just keep clicking
\textbf{Next}. A few worth knowing about:

\begin{itemize}[leftmargin=*]
    \item \textbf{Select Components}: default checkboxes are fine.
    \item \textbf{Choosing the default editor}: pick whichever text
    editor you're comfortable with (Notepad is a safe choice for
    beginners).
    \item \textbf{Adjusting your PATH environment}: choose
    ``Git from the command line and also from 3rd-party software''
    (this is usually the recommended/default option).
    \item \textbf{Choosing the SSH executable}: use the bundled
    OpenSSH (default option).
    \item \textbf{Configuring the line ending conversions}: keep the
    default (``Checkout Windows-style, commit Unix-style line
    endings'').
    \item \textbf{Configuring terminal emulator}: use MinTTY
    (default option).
\end{itemize}

\subsection{Finish Installation}
\begin{enumerate}[leftmargin=*]
    \item Click \textbf{Install} and wait for the process to complete.
    \item Click \textbf{Finish}.
\end{enumerate}

\subsection{Open Git Bash}
\begin{enumerate}[leftmargin=*]
    \item Right-click anywhere on your Desktop or inside a folder,
    and select \textbf{Git Bash Here} from the context menu.
    \item Alternatively, click the \textbf{Start} menu, search for
    ``Git Bash'', and open the application.
\end{enumerate}

\subsection{Verify the Installation}
In the Git Bash window that opens, type:

\begin{lstlisting}
git --version
\end{lstlisting}

You should see output similar to:

\begin{lstlisting}
git version 2.45.1.windows.1
\end{lstlisting}

If a version number appears, Git Bash is installed correctly and
you're ready to use commands like \texttt{git clone},
\texttt{ssh-keygen}, and \texttt{ssh -T git@github.com}.

\end{document}